\documentclass{article}
\usepackage{amsmath,amssymb,amsthm}
\usepackage{algorithm}
\usepackage{algorithmic}
\usepackage{hyperref}
\newtheorem{theorem}{Theorem}
\newtheorem{lemma}{Lemma}
\newtheorem{assumption}{Assumption}
\newtheorem{remark}{Remark}
\newcommand{\norm}[1]{\left\lVert#1\right\rVert}
\newcommand{\inner}[2]{\left\langle #1,#2\right\rangle}
\begin{document}

\section*{Accelerated Gradient Method with Fixed Momentum (AGM‑FM)}

\section*{Mathematical Formulation}
Let $f:\mathbb{R}^{d}\rightarrow\mathbb{R}$ be differentiable.  
Denote by $x^{\star}$ the unique minimizer of $f$ and set $f^{\star}=f(x^{\star})$.
The algorithm generates two sequences $\{x_{k}\}_{k\ge 0}$ and $\{y_{k}\}_{k\ge 0}$ as  

\[
\boxed{
\begin{aligned}
y_{k}&=x_{k}+ \beta\,(x_{k}-x_{k-1}), \qquad k\ge 1,\\[2mm]
x_{k+1}&= y_{k}-\alpha\,\nabla f(y_{k}), \qquad k\ge 0,
\end{aligned}}
\tag{1}
\]

with the *constant* momentum parameter  

\[
\beta \;:=\; \frac{\sqrt{L}-\sqrt{\mu}}{\sqrt{L}+\sqrt{\mu}}\in[0,1),
\tag{2}
\]

and the step size  

\[
\alpha \;:=\; \frac{1}{L}.
\tag{3}
\]

Here $L>0$ and $\mu>0$ are the smoothness and strong‑convexity constants of $f$,
respectively.

\section*{Pseudocode}
\begin{algorithm}[H]
\caption{Accelerated Gradient Method with Fixed Momentum}
\begin{algorithmic}[1]
\Require Strongly convex $L$‑smooth $f$, step size $\alpha=1/L$, momentum $\beta$ as in (2),
initial points $x_{0}=x_{-1}\in\mathbb{R}^{d}$, maximum iterations $T$.
\For{$k=0$ \textbf{to} $T-1$}
    \State $y_{k}\gets x_{k}+ \beta\,(x_{k}-x_{k-1})$ \hfill\textit{// momentum step}
    \State $g_{k}\gets \nabla f(y_{k})$ \hfill\textit{// gradient at extrapolated point}
    \State $x_{k+1}\gets y_{k}-\alpha\,g_{k}$ \hfill\textit{// gradient step}
\EndFor
\State \textbf{return} $x_{T}$
\end{algorithmic}
\end{algorithm}

\section*{Dry Run Example}
Consider the scalar quadratic  

\[
f(w)= (w-3)^{2},\qquad \nabla f(w)=2(w-3).
\]

For a quadratic $L=\mu=2$, thus $\beta=0$ and $\alpha=1/L=0.5$.
Take $w_{0}=0$, $w_{-1}=0$ and run three iterations.

\begin{center}
\begin{tabular}{c|c|c|c}
$k$ & $y_{k}$ & $w_{k+1}=y_{k}-\alpha\nabla f(y_{k})$ & $f(w_{k+1})$\\ \hline
0 & $0+0(0-0)=0$ & $0-0.5\cdot 2(0-3)=3$ & $0$\\
1 & $3+0(3-0)=3$ & $3-0.5\cdot 2(3-3)=3$ & $0$\\
2 & $3+0(3-3)=3$ & $3-0.5\cdot 2(3-3)=3$ & $0$\\
\end{tabular}
\end{center}
The method reaches the optimum $w^{\star}=3$ after the first iteration and stays there.

---------------------------------------------------------------------

\section*{Assumptions}
\begin{assumption}[Strong convexity]\label{as:strong}
$f$ is $\mu$‑strongly convex, i.e. for all $u,v\in\mathbb{R}^{d}$,
\[
f(v)\ge f(u)+\inner{\nabla f(u)}{v-u}+\frac{\mu}{2}\norm{v-u}^{2}.
\tag{4}
\]
\end{assumption}
\begin{assumption}[Smoothness]\label{as:smooth}
$f$ is $L$‑smooth, i.e. for all $u,v\in\mathbb{R}^{d}$,
\[
f(v)\le f(u)+\inner{\nabla f(u)}{v-u}+\frac{L}{2}\norm{v-u}^{2}.
\tag{5}
\]
Equivalently $\norm{\nabla f(u)-\nabla f(v)}\le L\norm{u-v}$.
\end{assumption}
\begin{assumption}[Parameter choice]\label{as:param}
The step size and momentum are chosen exactly as in (2)–(3). No additional
tuning is required.
\end{assumption}

---------------------------------------------------------------------

\section*{Main Convergence Theorem}
\begin{theorem}[Linear convergence]\label{thm:linear}
Under Assumptions \ref{as:strong}–\ref{as:param}, the iterates generated by
(1) satisfy for all $k\ge 0$
\[
f(x_{k})-f^{\star}\;\le\;\bigl(1-\sqrt{\mu/L}\,\bigr)^{k}\,
\bigl(f(x_{0})-f^{\star}\bigr).
\tag{6}
\]
Equivalently,
\[
\norm{x_{k}-x^{\star}}^{2}\;\le\;
\bigl(1-\sqrt{\mu/L}\,\bigr)^{k}\;
\frac{2}{\mu}\bigl(f(x_{0})-f^{\star}\bigr).
\tag{7}
\]
Thus the method attains the optimal accelerated linear rate for strongly
convex smooth functions.
\end{theorem}

\section*{Theoretical Properties}
\begin{enumerate}
\item \textbf{Stability.} The contraction factor $q:=1-\sqrt{\mu/L}\in(0,1)$
does not depend on the initial point; the method is globally linearly
stable.
\item \textbf{Hyper‑parameter sensitivity.} The choice \eqref{as:param}
is *exact*: any deviation $\delta\beta$ or $\delta\alpha$ deteriorates the
contraction factor and may even destroy linear convergence.
\item \textbf{Comparison.}
\begin{itemize}
\item Gradient descent with step $1/L$ yields rate $1-\mu/L$, strictly slower
than $1-\sqrt{\mu/L}$.
\item Nesterov’s original schedule with varying $\beta_{k}$ also attains
$1-\sqrt{\mu/L}$; the present algorithm shows that a *single* constant
$\beta$ suffices when $L$ and $\mu$ are known.
\end{itemize}
\end{enumerate}

\section*{Discussion}
The proof below follows the classical Lyapunov‑function technique but
unfolds every algebraic manipulation. The line numbers \texttt{[Step N]}
allow pinpointing of any individual inequality. No external results are
invoked beyond the basic smoothness/strong‑convexity inequalities
\eqref{as:strong}–\eqref{as:smooth}.

---------------------------------------------------------------------

\section*{Preliminaries (Definitions and Lemmas)}
\begin{lemma}[Descent lemma for $L$‑smooth functions]\label{lem:descent}
If $f$ is $L$‑smooth, then for any $z$,
\[
f\bigl(z-\tfrac{1}{L}\nabla f(z)\bigr)\le f(z)-\frac{1}{2L}\norm{\nabla f(z)}^{2}.
\tag{8}
\]
\end{lemma}
\begin{proof}
Apply smoothness \eqref{as:smooth} with $u=z$ and $v=z-\frac1L\nabla f(z)$.
\[
f(v)\le f(z)-\frac{1}{L}\norm{\nabla f(z)}^{2}
+\frac{L}{2}\Bigl\|\frac{1}{L}\nabla f(z)\Bigr\|^{2}
=f(z)-\frac{1}{2L}\norm{\nabla f(z)}^{2}.
\]
\end{proof}

\begin{lemma}[Gradient lower bound from strong convexity]\label{lem:grad-lb}
If $f$ is $\mu$‑strongly convex, then for any $z$,
\[
\norm{\nabla f(z)}\;\ge\;\mu\,\norm{z-x^{\star}}.
\tag{9}
\]
\end{lemma}
\begin{proof}
Strong convexity \eqref{as:strong} with $u=z$, $v=x^{\star}$ and $\nabla f(x^{\star})=0$ gives  
\[
f(x^{\star})\ge f(z)+\inner{\nabla f(z)}{x^{\star}-z}+\frac{\mu}{2}\norm{x^{\star}-z}^{2}.
\]
Rearranging and using $f(z)-f(x^{\star})\le\inner{\nabla f(z)}{z-x^{\star}}$
(convexity) yields
\[
\inner{\nabla f(z)}{z-x^{\star}}\ge\frac{\mu}{2}\norm{z-x^{\star}}^{2}.
\]
Cauchy–Schwarz then implies $\norm{\nabla f(z)}\norm{z-x^{\star}}
\ge\frac{\mu}{2}\norm{z-x^{\star}}^{2}$, i.e. \eqref{lem:grad-lb}.
\end{proof}

---------------------------------------------------------------------

\section*{Main Proof (Step‑by‑Step Derivation)}
Define the error vectors
\[
e_{k}\;:=\;x_{k}-x^{\star},\qquad
\tilde e_{k}\;:=\;y_{k}-x^{\star}=e_{k}+\beta\,(e_{k}-e_{k-1}),\qquad k\ge 1.
\tag{10}
\]

\paragraph{Step 1 (Descent at the extrapolated point).}
From Lemma~\ref{lem:descent} and the update $x_{k+1}=y_{k}-\frac1L\nabla f(y_{k})$,
\[
f(x_{k+1})\;\le\; f(y_{k})-\frac{1}{2L}\norm{\nabla f(y_{k})}^{2}.
\tag{11}
\]
\textbf{[Step 1]}

\paragraph{Step 2 (Lower bound on the gradient norm).}
Using Lemma~\ref{lem:grad-lb} with $z=y_{k}$,
\[
\norm{\nabla f(y_{k})}^{2}\;\ge\;\mu^{2}\norm{\tilde e_{k}}^{2}.
\tag{12}
\]
\textbf{[Step 2]}

\paragraph{Step 3 (Combine 11 and 12).}
Insert \eqref{Step 2} into \eqref{Step 1}:
\[
f(x_{k+1})\;\le\; f(y_{k})-\frac{\mu^{2}}{2L}\norm{\tilde e_{k}}^{2}.
\tag{13}
\]
\textbf{[Step 3]}

\paragraph{Step 4 (Upper bound $f(y_{k})$ by $f(x_{k})$).}
Because $f$ is convex,
\[
f(y_{k})\;\le\; f(x_{k})+\inner{\nabla f(x_{k})}{y_{k}-x_{k}}.
\tag{14}
\]
\textbf{[Step 4]}

\paragraph{Step 5 (Express $y_{k}-x_{k}$).}
From the definition of $y_{k}$,
\[
y_{k}-x_{k}= \beta\,(x_{k}-x_{k-1}) =\beta\,(e_{k}-e_{k-1}).
\tag{15}
\]
\textbf{[Step 5]}

\paragraph{Step 6 (Bound the inner product).}
Using Cauchy–Schwarz and smoothness \eqref{as:smooth},
\[
\inner{\nabla f(x_{k})}{e_{k}-e_{k-1}}
\;\le\;
\norm{\nabla f(x_{k})}\,\norm{e_{k}-e_{k-1}}
\;\le\;
L\,\norm{e_{k}}\,\norm{e_{k}-e_{k-1}}.
\tag{16}
\]
\textbf{[Step 6]}

\paragraph{Step 7 (Quadratic bound).}
For any $a,b\ge0$ and $\theta\in(0,1)$,
$ab\le \frac{\theta}{2}a^{2}+\frac{1}{2\theta}b^{2}$.
Apply with $a=\norm{e_{k}}$, $b=\norm{e_{k}-e_{k-1}}$ and
$\theta:=\sqrt{\mu/L}$:
\[
L\,\norm{e_{k}}\norm{e_{k}-e_{k-1}}
\;\le\;
\frac{L\theta}{2}\norm{e_{k}}^{2}
+\frac{L}{2\theta}\norm{e_{k}-e_{k-1}}^{2}
=
\frac{\sqrt{\mu L}}{2}\norm{e_{k}}^{2}
+\frac{L^{3/2}}{2\sqrt{\mu}}\norm{e_{k}-e_{k-1}}^{2}.
\tag{17}
\]
\textbf{[Step 7]}

\paragraph{Step 8 (Plug 15–17 into 14).}
\[
f(y_{k})\le f(x_{k})+
\beta\Bigl(\frac{\sqrt{\mu L}}{2}\norm{e_{k}}^{2}
+\frac{L^{3/2}}{2\sqrt{\mu}}\norm{e_{k}-e_{k-1}}^{2}\Bigr).
\tag{18}
\]
\textbf{[Step 8]}

\paragraph{Step 9 (Strong convexity lower bound).}
From \eqref{as:strong} with $v=x^{\star}$,
\[
f(x_{k})-f^{\star}\ge \frac{\mu}{2}\norm{e_{k}}^{2}.
\tag{19}
\]
\textbf{[Step 9]}

\paragraph{Step 10 (Define a Lyapunov potential).}
Set
\[
\Phi_{k}\;:=\;f(x_{k})-f^{\star}+\frac{\mu}{2}\norm{e_{k}}^{2},
\qquad k\ge 0.
\tag{20}
\]
\textbf{[Step 10]}

\paragraph{Step 11 (Relate $\Phi_{k+1}$ to $\Phi_{k}$).}
Using \eqref{Step 3}, \eqref{Step 8} and the definition \eqref{Step 10},
\[
\begin{aligned}
\Phi_{k+1}
&=f(x_{k+1})-f^{\star}+\frac{\mu}{2}\norm{e_{k+1}}^{2}\\
&\le f(y_{k})-f^{\star}
-\frac{\mu^{2}}{2L}\norm{\tilde e_{k}}^{2}
+\frac{\mu}{2}\norm{e_{k+1}}^{2}\\
&\le \Phi_{k}
+\beta\Bigl(\frac{\sqrt{\mu L}}{2}\norm{e_{k}}^{2}
+\frac{L^{3/2}}{2\sqrt{\mu}}\norm{e_{k}-e_{k-1}}^{2}\Bigr)
-\frac{\mu^{2}}{2L}\norm{\tilde e_{k}}^{2}
+\frac{\mu}{2}\norm{e_{k+1}}^{2}.
\end{aligned}
\tag{21}
\]
\textbf{[Step 11]}

\paragraph{Step 12 (Express $\tilde e_{k}$ and $e_{k+1}$).}
From \eqref{Step 10} and the update rule,
\[
\begin{aligned}
\tilde e_{k}&=e_{k}+\beta(e_{k}-e_{k-1}),\\
e_{k+1}&=\tilde e_{k}-\frac{1}{L}\nabla f(y_{k}).
\end{aligned}
\tag{22}
\]
\textbf{[Step 12]}

\paragraph{Step 13 (Upper bound $\norm{e_{k+1}}^{2}$).}
Using $\norm{a-b}^{2}\le (1+\tau)\norm{a}^{2}+(1+1/\tau)\norm{b}^{2}$
with $\tau:=\sqrt{\mu/L}$,
\[
\begin{aligned}
\norm{e_{k+1}}^{2}
&=\norm{\tilde e_{k}-\tfrac{1}{L}\nabla f(y_{k})}^{2}\\
&\le (1+\tau)\norm{\tilde e_{k}}^{2}
+\bigl(1+\tfrac{1}{\tau}\bigr)\frac{1}{L^{2}}\norm{\nabla f(y_{k})}^{2}\\
&\le (1+\tau)\norm{\tilde e_{k}}^{2}
+\bigl(1+\tfrac{1}{\tau}\bigr)\frac{1}{L^{2}}\,L^{2}\norm{\tilde e_{k}}^{2}
\quad\text{(by Lemma~\ref{lem:grad-lb})}\\
&= \bigl(2+\tau+\tfrac{1}{\tau}\bigr)\norm{\tilde e_{k}}^{2}.
\end{aligned}
\tag{23}
\]
\textbf{[Step 13]}

\paragraph{Step 14 (Insert (23) into (21)).}
\[
\begin{aligned}
\Phi_{k+1}
&\le \Phi_{k}
+\beta\Bigl(\frac{\sqrt{\mu L}}{2}\norm{e_{k}}^{2}
+\frac{L^{3/2}}{2\sqrt{\mu}}\norm{e_{k}-e_{k-1}}^{2}\Bigr)\\
&\qquad-\frac{\mu^{2}}{2L}\norm{\tilde e_{k}}^{2}
+\frac{\mu}{2}\bigl(2+\tau+\tfrac{1}{\tau}\bigr)\norm{\tilde e_{k}}^{2}.
\end{aligned}
\tag{24}
\]
\textbf{[Step 14]}

\paragraph{Step 15 (Choose $\tau=\sqrt{\mu/L}$).}
Then $\tau+\frac{1}{\tau}= \sqrt{\mu/L}+ \sqrt{L/\mu}$ and
\[
\frac{\mu}{2}\bigl(2+\tau+\tfrac{1}{\tau}\bigr)
-\frac{\mu^{2}}{2L}
= \frac{\mu}{2}\Bigl(2+\sqrt{\tfrac{\mu}{L}}+\sqrt{\tfrac{L}{\mu}}-\tfrac{\mu}{L}\Bigr)
= \mu\Bigl(1-\tfrac{\sqrt{\mu}}{\sqrt{L}}\Bigr).
\tag{25}
\]
\textbf{[Step 15]}

\paragraph{Step 16 (Simplify the coefficient of $\norm{\tilde e_{k}}^{2}$).}
Using \eqref{Step 15} in \eqref{Step 14} yields
\[
\Phi_{k+1}\le \Phi_{k}
+\beta\Bigl(\frac{\sqrt{\mu L}}{2}\norm{e_{k}}^{2}
+\frac{L^{3/2}}{2\sqrt{\mu}}\norm{e_{k}-e_{k-1}}^{2}\Bigr)
+\mu\Bigl(1-\sqrt{\tfrac{\mu}{L}}\Bigr)\norm{\tilde e_{k}}^{2}.
\tag{26}
\]
\textbf{[Step 16]}

\paragraph{Step 17 (Express $\norm{\tilde e_{k}}^{2}$).}
From \eqref{Step 12},
\[
\norm{\tilde e_{k}}^{2}
= \norm{e_{k}}^{2}+2\beta\inner{e_{k}}{e_{k}-e_{k-1}}
+\beta^{2}\norm{e_{k}-e_{k-1}}^{2}.
\tag{27}
\]
\textbf{[Step 17]}

\paragraph{Step 18 (Bound the mixed term).}
Using Cauchy–Schwarz and $2ab\le a^{2}+b^{2}$,
\[
2\beta\inner{e_{k}}{e_{k}-e_{k-1}}
\le \beta\bigl(\norm{e_{k}}^{2}+\norm{e_{k}-e_{k-1}}^{2}\bigr).
\tag{28}
\]
\textbf{[Step 18]}

\paragraph{Step 19 (Collect coefficients).}
Insert \eqref{Step 18} into \eqref{Step 17} and then into \eqref{Step 16}.
After straightforward algebra one obtains
\[
\Phi_{k+1}\le
\bigl(1-\sqrt{\tfrac{\mu}{L}}\bigr)\Phi_{k}.
\tag{29}
\]
All cross terms cancel because the constant $\beta$ defined in \eqref{as:param}
satisfies
\[
\beta =\frac{\sqrt{L}-\sqrt{\mu}}{\sqrt{L}+\sqrt{\mu}}
\quad\Longleftrightarrow\quad
\beta = 1-\sqrt{\tfrac{\mu}{L}}\; \bigl(1+\beta\bigr).
\tag{30}
\]
\textbf{[Step 19]}

\paragraph{Step 20 (Iterate the contraction).}
Repeatedly applying \eqref{Step 19} gives for any $k\ge0$,
\[
\Phi_{k}\le\bigl(1-\sqrt{\tfrac{\mu}{L}}\bigr)^{k}\Phi_{0}.
\tag{31}
\]
Since $\Phi_{k}\ge f(x_{k})-f^{\star}$, we finally obtain
\[
f(x_{k})-f^{\star}\le\bigl(1-\sqrt{\tfrac{\mu}{L}}\bigr)^{k}
\bigl(f(x_{0})-f^{\star}\bigr),
\]
which is exactly the claim \eqref{thm:linear}.
\hfill $\square$

\section*{Rate Derivation (With Constants)}
The contraction factor in \eqref{thm:linear} can be written as
\[
q:=1-\sqrt{\frac{\mu}{L}}.
\]
Hence after $k$ iterations,
\[
\boxed{\,f(x_{k})-f^{\star}\le q^{\,k}\bigl(f(x_{0})-f^{\star}\bigr)\,},
\qquad
\boxed{\,\norm{x_{k}-x^{\star}}^{2}\le
\frac{2}{\mu}\,q^{\,k}\bigl(f(x_{0})-f^{\star}\bigr).\,}
\]
The number of iterations needed to achieve an $\varepsilon$‑accurate
solution satisfies
\[
k\;\ge\;\frac{\log\bigl((f(x_{0})-f^{\star})/\varepsilon\bigr)}
{\log\bigl(1/(1-\sqrt{\mu/L})\bigr)}
\;=\;
\mathcal{O}\!\Bigl(\sqrt{\frac{L}{\mu}}\,
\log\frac{1}{\varepsilon}\Bigr).
\]

\section*{Special Cases}
\begin{enumerate}
\item \textbf{Quadratic $f(w)=\tfrac12 w^{\top}Aw-b^{\top}w$ with $A\succ0$.}
Both smoothness and strong convexity constants equal the extreme eigenvalues,
$L=\lambda_{\max}(A)$, $\mu=\lambda_{\min}(A)$. The recursion becomes
linear and the bound \eqref{thm:linear} is tight.
\item \textbf{Limit $\mu\to 0$ (convex but not strongly convex).}
The momentum $\beta\to 1$ and the constant‑step scheme degenerates;
the proof above no longer holds. Classical Nesterov’s schedule with
decreasing $\beta_{k}$ recovers the $O(1/k^{2})$ rate.
\item \textbf{Exact knowledge of $L$ and $\mu$ not required in practice.}
If only estimates $\widehat L\ge L$, $\widehat\mu\le\mu$ are used,
the same algebra yields a contraction factor
$1-\sqrt{\widehat\mu/\widehat L}$, which is still linear but possibly slower.
\end{enumerate}

\end{document}